\documentclass{zjureport}
\special{dvipdfmx:config z 0} % 取消PDF压缩，加快速度，最终版本生成的时候最好把这句话注释掉

% =============================================
% Part 0 Edit the info
% =============================================

\major{信息工程}
\name{黄鹤翔}
\partner{}
\title{本科实验报告}
\stuid{3230106231}
\college{信息与电子工程学院}
\date{2025-3-26}
\lab{东四224}
\course{场波实验}
\instructor{王子立}
\grades{}
\expname{微带传输线负载特性分析}
\exptype{探究型}

\usepackage{listings}
\usepackage{xcolor}

\lstset{
 columns=fixed,       
 numbers=left,                                        % 在左侧显示行号
 numberstyle=\tiny\color{gray},                       % 设定行号格式
 frame=none,                                          % 不显示背景边框
 backgroundcolor=\color[RGB]{245,245,244},            % 设定背景颜色
 keywordstyle=\color[RGB]{40,40,255},                 % 设定关键字颜色
 numberstyle=\footnotesize\color{darkgray},           
 commentstyle=\it\color[RGB]{0,96,96},                % 设置代码注释的格式
 stringstyle=\rmfamily\slshape\color[RGB]{128,0,0},   % 设置字符串格式
 showstringspaces=false,                              % 不显示字符串中的空格
 breaklines=true,                                     % 对过长的代码自动换行
 language=verilog,                                     % 设置语言
}


\begin{document}
% =============================================
% Part 1 Header
% =============================================
\makecover

% \makecontent

\makeheader
% =============================================
% Part 2 Main document
% =============================================

% \begin{figure}[H] %H为当前位置，!htb为忽略美学标准，htbp为浮动图形
%     \centering %图片居中
%     \includegraphics[width=0.8\textwidth]{fig1-SR830structure.jpg}
%     \caption{锁相放大器SR830结构图} %最终文档中希望显示的图片标题
% \end{figure}

\section{实验目的}
\begin{itemize}
  \item 了解基本传输线、微带线的特性
  \item 熟悉网络参量测量，掌握矢网网络分析仪的基本使用方法
\end{itemize}

\section{实验原理}
\begin{figure}[H]
  \centering
  \includegraphics[width=0.8\textwidth]{原理图.png}
  \caption{微带线负载特性矢网网络分析原理图} 
\end{figure}

考虑一段特性阻抗为$Z_0$的传输线，一端接信号源，另一端则接上负载。
假设此传输线无耗，且传输线$\gamma = j \beta$，则传输线上电压及电流可用
下列二式表示：
\begin{align*}
  U(z) = U^+e^{-\beta z} + U^-e^{\beta z} \\
  I(z) = I^+e^{-\beta z} - I^-e^{\beta z}
\end{align*}

\subsection{负载端(z = 0)处情况}
电压及电流为\par
\begin{align*}
  U = U_L = U^+ + U^- \\
  I = I_L = I^+ - I^-
\end{align*}
而$Z_0I^+ = U^+,Z_0I^- = U^-$，公式可改写成\par
\begin{align*}
  I_L = \frac{1}{Z_0}(U^+-U^-)
\end{align*}
可得负载阻抗为\par
\begin{align*}
  Z_L = \frac{U_L}{I_L} = Z_0\frac{U^+ + U^-}{U^+ - U^-}
\end{align*}
定义归一化负载阻抗为\par
\begin{align*}
  z_L = Z_L = \frac{Z_L}{Z_0} = \frac{1+\Gamma_L}{1-\Gamma_L}
\end{align*}
其中定义$\Gamma_L$为负载端的电压反射系数\par
\begin{align*}
  \Gamma_L = \frac{U^-}{U^+} = \frac{Z_L - Z_0}{Z_L + Z_0} = |\Gamma_L|e^{j\phi_L}
\end{align*}
当$Z_L = Z_0$或为无限长传输线时，$\Gamma_L = 0$，无反射波，是行波状态或匹配状态。\par
当$Z_L$为纯电抗元件或处于开路或者短路状态时，$|\Gamma_L| = 1$，全反射，是驻波状态。\par
当$Z_L$为其他值时，$|\Gamma_L| <= 1$，为行波驻波状态。\par
线上任意点的反射系数为\par
\begin{align*}
  \Gamma_L = |\Gamma_L|e^{j\phi_L-j2\beta z}
\end{align*}
定义驻波比VSWR和回波损耗RL为\par
\begin{align*}
  VSWR = \frac{1+|\Gamma_L|}{1-|\Gamma_L|}\\
  RL = -20\log_{10}|\Gamma_L|
\end{align*}

\subsection{输入端(z = -l)处情况}
反射系数$\Gamma(z)$应改成\par
\begin{align*}
  \Gamma(L) = \frac{U^-e^{-j\beta_L}}{U^+e^{j\phi\beta_L}} = \frac{U^-}{U^+}e^{-j2\beta_L} = \Gamma_L e^{-j2\beta_L}
\end{align*}
输入阻抗为\par
\begin{align*}
  Z_{in} = Z_0\frac{Z_L+jZ_0\tan(\beta L)}{Z_0+jZ_L\tan(\beta L)}
\end{align*}
由上式可知\par
\begin{itemize}
  \item 当$L \rightarrow \infty$时，$Z_{in} \rightarrow Z_0$
  \item 当$L = \frac{\lambda}{2}$时，$Z_{in} = Z_L$
  \item 当$L = \frac{\lambda}{4}$时，$Z_{in} = \frac{Z_0^2}{Z_L}$
\end{itemize}

\section{实验内容}
\begin{itemize}
  \item 了解矢量网络分析仪的原理和使用方法
  \item 用矢量网络分析仪分别测量微带开路传输线模块的反射特性，
  并引入电阻、电容和电感负载测量并分析在不同负载情况下的反射特性
\end{itemize}

\section{实验过程与记录}
\subsection{微带开路传输线模块}
\subsubsection{校准}
连接好射频电缆线后，按【校准】键，选择[机械校准],[单端口(反射)]，
分别接上开路器、断路器、匹配负载，最后按下[完成单端口]键，完成校准，校准
后接匹配负载的史密斯圆图如下图所示：\par
\begin{figure}[H]
  \centering
  \includegraphics[width=0.8\textwidth,angle = 180]{微带传输线校准.jpg}
  \caption{校准后的史密斯圆图}
\end{figure}

\subsubsection{模块上不加负载}
\begin{figure}[H]
  \centering
  \includegraphics[width=0.8\textwidth,angle = 180]{传输线模块开路.jpg}
  \caption{模块开路}
\end{figure}
移动光标至2.5GHz处，光标接近但略微偏离开路点，可能
是因为微带线的开路段并非理想开路，实际存在边缘场效应，
导致等效电长度略长于物理长度，引起误差。\par

\subsubsection{模块接匹配电阻}
\begin{figure}[H]
  \centering
  \includegraphics[width=0.8\textwidth,angle = 180]{传输线模块49.9电阻.jpg}
  \caption{模块匹配}
\end{figure}
曲线在匹配点附近\par

\subsubsection{模块接0$\Omega$电阻}
\begin{figure}[H]
  \centering
  \includegraphics[width=0.8\textwidth,angle = 180]{传输线模块短路.jpg}
  \caption{模块接0$\Omega$电阻}  
\end{figure}
理论上短路时曲线应该在左端点附近，但实验测得曲线在右端点附近，
可能是在接入0$\Omega$电阻时，没有让电阻与模块完全接触，
造成了类似断路的现象，导致实验值与理论值相差很大。\par

\subsubsection{模块接电容}
\begin{figure}[H]
  \centering
  \includegraphics[width=0.8\textwidth,angle = 180]{传输线模块电容.jpg}
  \caption{模块接电容}
\end{figure}
可以看到曲线大部分落在下半圆内，呈现容性特征。

\subsubsection{模块接电感}
\begin{figure}[H]
  \centering
  \includegraphics[width=0.8\textwidth,angle = 180]{传输线模块电感.jpg}
  \caption{模块接电感}
\end{figure}
可以看到曲线大部分落在上半圆内，呈现感性特征。

\subsection{天线}
取下射频电缆线，重新校准后接上天线，得到的史密斯圆图如下图所示：\par
\begin{figure}[H]
  \centering
  \includegraphics[width=0.8\textwidth,angle = 180]{天线.jpg}
  \caption{天线测量的史密斯圆图}
\end{figure}

\subsection{微带耦合滤波器}

\subsubsection{基础参数测量}
中心频率为2.3675GHz，3dB带宽为0.161GHz，带内波纹约为1dB，
阻带衰减36.484dB\par
\begin{figure}[H]
  \centering
  \includegraphics[width=0.8\textwidth,angle = 180]{中心频率与带宽.jpg}
\end{figure}
观察光标2可得中心频率为2.3675GHz，观察光标6、7可得3dB带宽为0.161GHz，
可大致估计带内波纹为0.2dB\par
\begin{figure}[H]
  \centering
  \includegraphics[width=0.8\textwidth,angle = 180]{阻带衰减.jpg}
\end{figure}
观察光标6可得阻带衰减为36.484dB。

\subsubsection{S参数与相位}
\begin{figure}[H]
  \centering
  \includegraphics[width=0.8\textwidth,angle = 180]{无干扰swr.jpg}
  \caption{$S_{21}$}
\end{figure}
\begin{figure}[H]
  \centering
  \includegraphics[width=0.8\textwidth,angle = 180]{相位.jpg}
  \caption{相位}
\end{figure}

\subsubsection{受到干扰的变化}
\begin{figure}[H]
  \begin{minipage}{0.45\textwidth}
    \centering
    \includegraphics[width=0.8\textwidth,angle = 180]{无干扰特性.jpg}
    \caption{干扰前特性曲线}
  \end{minipage}
  \begin{minipage}{0.45\textwidth}
    \centering
    \includegraphics[width=0.8\textwidth,angle = 180]{干扰特性.jpg}
    \caption{干扰后特性曲线}
  \end{minipage}
  \caption{特性曲线变化}
\end{figure}

\begin{figure}[H]
  \begin{minipage}{0.45\textwidth}
    \centering
    \includegraphics[width=0.8\textwidth,angle = 180]{无干扰swr.jpg}
    \caption{干扰前驻波比}
  \end{minipage}
  \begin{minipage}{0.45\textwidth}
    \centering
    \includegraphics[width=0.8\textwidth,angle = 180]{干扰swr.jpg}
    \caption{干扰后驻波比}
  \end{minipage}
  \caption{驻波比变化}
\end{figure}

\begin{figure}[H]
  \begin{minipage}{0.45\textwidth}
    \centering
    \includegraphics[width=0.8\textwidth,angle = 180]{无干扰圆图.jpg}
    \caption{干扰前圆图}
  \end{minipage}
  \begin{minipage}{0.45\textwidth}
    \centering
    \includegraphics[width=0.8\textwidth,angle = 180]{干扰圆图.jpg}
    \caption{干扰后圆图}
  \end{minipage}
  \caption{圆图变化}
\end{figure}

由上述数据可知，微带滤波器为带通滤波器，中心频率2.3675GHz，3dB截止
频率为2.448GHz和2.287GHz，带宽仅为0.161GHz，带宽较小。\\
由$|S_{21}|$与$\omega$的关系图可以看到，
滤波器在很大一部分$\omega$内反射系数均高于1.5，反射损耗较大；
由$S_{21}$的相位图可知，滤波器的线性性较好。\\
综上所述，微带滤波器的性能一般。

\subsubsection{输入阻抗与反射系数}
\begin{figure}[H]
  \centering
  \includegraphics[width=0.8\textwidth,angle = 180]{无干扰圆图.jpg}
\end{figure}
\end{document}